
%%%%%%%%%%%%%%%%%%%%%%%%%%%%%%%%%%%%%%%%%%%%%%%%%
%%  Section: Matrix Functions
%%%%%%%%%%%%%%%%%%%%%%%%%%%%%%%%%%%%%%%%%%%%%%%%%
\section{Matrix Functions}

%%%%%%%%%%%%%%%%%%%%%%
%%%     Slide      %%%
%%%%%%%%%%%%%%%%%%%%%% 
\begin{frame}[t]{Matrix Functions}
    For a diagonalizable matrix $A = T \Lambda T^{-1}$, we can extend a scalar function $f$ (such as polynomials, exponentials, or roots) to the matrix case.

    \begin{alert}
        If $f$ is defined on the spectrum%%%%%%%%%%%%
        \footnote{The \emph{spectrum} of a matrix $A$, denoted $\sigma(A)$, is the set of all eigenvalues of $A$, i.e.,
            \[
                \sigma(A) = \{\lambda \in \mathbb{C} : \det(A - \lambda I) = 0\}.
            \]}%
        of $A$, we define the \textbf{matrix function} as:
        \[
            f(A) = T \, f(\Lambda) \, T^{-1}
        \]
    \end{alert}


    where $f(\Lambda)$ is the diagonal matrix obtained by applying $f$ to each eigenvalue:
    \[
        f(\Lambda) = \operatorname{diag}\left(f(\lambda_1), f(\lambda_2), \dots, f(\lambda_n)\right)
    \]


    \vfill
\end{frame}


%%%%%%%%%%%%%%%%%%%%%%
%%%     Slide      %%%
%%%%%%%%%%%%%%%%%%%%%% 
\begin{frame}[t,allowframebreaks]{Matrix Powers --}

    \textbf{Why diagonalize?}
    Computing powers $A^n$ directly requires $n$ matrix multiplications.
    But if $A$ is diagonalizable, we can write
    \[
        A = T \Lambda T^{-1},
    \]
    and therefore
    \[
        A^n = (T \Lambda T^{-1})^n = T\, \Lambda^n\, T^{-1}.
    \]

    Since $\Lambda$ is diagonal, its powers are trivial:
    \[
        \Lambda^n =
        \mathrm{diag}(\lambda_1^n, \lambda_2^n, \dots, \lambda_n^n),
    \]
    so the computation becomes extremely inexpensive.

    \medskip
    \pagebreak

    \textbf{Examples:}
    \[
        A^{1/2} = T\, \Lambda^{1/2}\, T^{-1},
        \qquad
        e^A = T\, e^{\Lambda}\, T^{-1},
    \]
    with
    \[
        e^\Lambda =
        \mathrm{diag}(e^{\lambda_1},\, \dots,\, e^{\lambda_n}).
    \]

    \begin{alertblock}{Why this matters (especially in EDA):}\smallskip
        Diagonalization converts difficult matrix computations into simple scalar
        operations on eigenvalues.
        This enables fast evaluation of $A^n$, $A^{1/2}$, and $e^A$, which arise in
        simulation algorithms, stability analysis, and time‑domain state‑transition solutions.
    \end{alertblock}

    \vfill

\end{frame}

\begin{frame}[t, allowframebreaks]{Matrix Exponential --}

    \textbf{Scalar exponential (review):}
    \[
        e^{x} = 1 + x + \frac{x^{2}}{2!} + \frac{x^{3}}{3!} + \cdots
    \]

    \medskip

    \textbf{Matrix exponential:}
    For any square matrix $A$, the exponential is defined by the same power series:
    \[
        e^{A} \;=\; \sum_{n=0}^{\infty} \frac{A^{n}}{n!}.
    \]

    \medskip

    \textbf{Key idea:}
    The series converges for every matrix $A$, and $e^{A}$ behaves like a
    generalization of the scalar exponential to matrices.

    \medskip


    \textbf{1. Power–series definition:}
    \[
        e^{A}
        = \sum_{n=0}^{\infty} \frac{A^{n}}{n!}.
    \]

    \medskip

    \textbf{2. Diagonalization:}
    If $A$ is diagonalizable,
    \[
        A = T \Lambda T^{-1},
    \]
    where $\Lambda = \mathrm{diag}(\lambda_1,\dots,\lambda_n)$.

    \medskip

    \textbf{3. Put the series into the diagonal form:}
    \[
        A^{n} = (T \Lambda T^{-1})^{n}
        = T\,\Lambda^{n}\,T^{-1}.
    \]

    \medskip

    \textbf{4. Substitute into the series:}
    \[
        e^{A}
        = \sum_{n=0}^{\infty} \frac{A^{n}}{n!}
        = \sum_{n=0}^{\infty}
        \frac{T \Lambda^{n} T^{-1}}{n!}
        = T\left(\sum_{n=0}^{\infty} \frac{\Lambda^{n}}{n!}\right)T^{-1}.
    \]

    \medskip

    \textbf{5. But $\Lambda$ is diagonal, so}
    \[
        e^{\Lambda}
        = \mathrm{diag}(e^{\lambda_1},\, \dots,\, e^{\lambda_n}),
    \]
    and therefore
    \[
        \boxed{\,e^{A} = T e^{\Lambda} T^{-1}\, }.
    \]

\end{frame}


%%EoF